\documentclass[11pt]{article}

    \usepackage[breakable]{tcolorbox}
    \usepackage{parskip} % Stop auto-indenting (to mimic markdown behaviour)
    
    \usepackage{iftex}
    \ifPDFTeX
    	\usepackage[T1]{fontenc}
    	\usepackage{mathpazo}
    \else
    	\usepackage{fontspec}
    \fi

    % Basic figure setup, for now with no caption control since it's done
    % automatically by Pandoc (which extracts ![](path) syntax from Markdown).
    \usepackage{graphicx}
    % Maintain compatibility with old templates. Remove in nbconvert 6.0
    \let\Oldincludegraphics\includegraphics
    % Ensure that by default, figures have no caption (until we provide a
    % proper Figure object with a Caption API and a way to capture that
    % in the conversion process - todo).
    \usepackage{caption}
    \DeclareCaptionFormat{nocaption}{}
    \captionsetup{format=nocaption,aboveskip=0pt,belowskip=0pt}

    \usepackage[Export]{adjustbox} % Used to constrain images to a maximum size
    \adjustboxset{max size={0.9\linewidth}{0.9\paperheight}}
    \usepackage{float}
    \floatplacement{figure}{H} % forces figures to be placed at the correct location
    \usepackage{xcolor} % Allow colors to be defined
    \usepackage{enumerate} % Needed for markdown enumerations to work
    \usepackage{geometry} % Used to adjust the document margins
    \usepackage{amsmath} % Equations
    \usepackage{amssymb} % Equations
    \usepackage{textcomp} % defines textquotesingle
    % Hack from http://tex.stackexchange.com/a/47451/13684:
    \AtBeginDocument{%
        \def\PYZsq{\textquotesingle}% Upright quotes in Pygmentized code
    }
    \usepackage{upquote} % Upright quotes for verbatim code
    \usepackage{eurosym} % defines \euro
    \usepackage[mathletters]{ucs} % Extended unicode (utf-8) support
    \usepackage{fancyvrb} % verbatim replacement that allows latex
    \usepackage{grffile} % extends the file name processing of package graphics 
                         % to support a larger range
    \makeatletter % fix for grffile with XeLaTeX
    \def\Gread@@xetex#1{%
      \IfFileExists{"\Gin@base".bb}%
      {\Gread@eps{\Gin@base.bb}}%
      {\Gread@@xetex@aux#1}%
    }
    \makeatother

    % The hyperref package gives us a pdf with properly built
    % internal navigation ('pdf bookmarks' for the table of contents,
    % internal cross-reference links, web links for URLs, etc.)
    \usepackage{hyperref}
    % The default LaTeX title has an obnoxious amount of whitespace. By default,
    % titling removes some of it. It also provides customization options.
    \usepackage{titling}
    \usepackage{longtable} % longtable support required by pandoc >1.10
    \usepackage{booktabs}  % table support for pandoc > 1.12.2
    \usepackage[inline]{enumitem} % IRkernel/repr support (it uses the enumerate* environment)
    \usepackage[normalem]{ulem} % ulem is needed to support strikethroughs (\sout)
                                % normalem makes italics be italics, not underlines
    \usepackage{mathrsfs}
    

    
    % Colors for the hyperref package
    \definecolor{urlcolor}{rgb}{0,.145,.698}
    \definecolor{linkcolor}{rgb}{.71,0.21,0.01}
    \definecolor{citecolor}{rgb}{.12,.54,.11}

    % ANSI colors
    \definecolor{ansi-black}{HTML}{3E424D}
    \definecolor{ansi-black-intense}{HTML}{282C36}
    \definecolor{ansi-red}{HTML}{E75C58}
    \definecolor{ansi-red-intense}{HTML}{B22B31}
    \definecolor{ansi-green}{HTML}{00A250}
    \definecolor{ansi-green-intense}{HTML}{007427}
    \definecolor{ansi-yellow}{HTML}{DDB62B}
    \definecolor{ansi-yellow-intense}{HTML}{B27D12}
    \definecolor{ansi-blue}{HTML}{208FFB}
    \definecolor{ansi-blue-intense}{HTML}{0065CA}
    \definecolor{ansi-magenta}{HTML}{D160C4}
    \definecolor{ansi-magenta-intense}{HTML}{A03196}
    \definecolor{ansi-cyan}{HTML}{60C6C8}
    \definecolor{ansi-cyan-intense}{HTML}{258F8F}
    \definecolor{ansi-white}{HTML}{C5C1B4}
    \definecolor{ansi-white-intense}{HTML}{A1A6B2}
    \definecolor{ansi-default-inverse-fg}{HTML}{FFFFFF}
    \definecolor{ansi-default-inverse-bg}{HTML}{000000}

    % commands and environments needed by pandoc snippets
    % extracted from the output of `pandoc -s`
    \providecommand{\tightlist}{%
      \setlength{\itemsep}{0pt}\setlength{\parskip}{0pt}}
    \DefineVerbatimEnvironment{Highlighting}{Verbatim}{commandchars=\\\{\}}
    % Add ',fontsize=\small' for more characters per line
    \newenvironment{Shaded}{}{}
    \newcommand{\KeywordTok}[1]{\textcolor[rgb]{0.00,0.44,0.13}{\textbf{{#1}}}}
    \newcommand{\DataTypeTok}[1]{\textcolor[rgb]{0.56,0.13,0.00}{{#1}}}
    \newcommand{\DecValTok}[1]{\textcolor[rgb]{0.25,0.63,0.44}{{#1}}}
    \newcommand{\BaseNTok}[1]{\textcolor[rgb]{0.25,0.63,0.44}{{#1}}}
    \newcommand{\FloatTok}[1]{\textcolor[rgb]{0.25,0.63,0.44}{{#1}}}
    \newcommand{\CharTok}[1]{\textcolor[rgb]{0.25,0.44,0.63}{{#1}}}
    \newcommand{\StringTok}[1]{\textcolor[rgb]{0.25,0.44,0.63}{{#1}}}
    \newcommand{\CommentTok}[1]{\textcolor[rgb]{0.38,0.63,0.69}{\textit{{#1}}}}
    \newcommand{\OtherTok}[1]{\textcolor[rgb]{0.00,0.44,0.13}{{#1}}}
    \newcommand{\AlertTok}[1]{\textcolor[rgb]{1.00,0.00,0.00}{\textbf{{#1}}}}
    \newcommand{\FunctionTok}[1]{\textcolor[rgb]{0.02,0.16,0.49}{{#1}}}
    \newcommand{\RegionMarkerTok}[1]{{#1}}
    \newcommand{\ErrorTok}[1]{\textcolor[rgb]{1.00,0.00,0.00}{\textbf{{#1}}}}
    \newcommand{\NormalTok}[1]{{#1}}
    
    % Additional commands for more recent versions of Pandoc
    \newcommand{\ConstantTok}[1]{\textcolor[rgb]{0.53,0.00,0.00}{{#1}}}
    \newcommand{\SpecialCharTok}[1]{\textcolor[rgb]{0.25,0.44,0.63}{{#1}}}
    \newcommand{\VerbatimStringTok}[1]{\textcolor[rgb]{0.25,0.44,0.63}{{#1}}}
    \newcommand{\SpecialStringTok}[1]{\textcolor[rgb]{0.73,0.40,0.53}{{#1}}}
    \newcommand{\ImportTok}[1]{{#1}}
    \newcommand{\DocumentationTok}[1]{\textcolor[rgb]{0.73,0.13,0.13}{\textit{{#1}}}}
    \newcommand{\AnnotationTok}[1]{\textcolor[rgb]{0.38,0.63,0.69}{\textbf{\textit{{#1}}}}}
    \newcommand{\CommentVarTok}[1]{\textcolor[rgb]{0.38,0.63,0.69}{\textbf{\textit{{#1}}}}}
    \newcommand{\VariableTok}[1]{\textcolor[rgb]{0.10,0.09,0.49}{{#1}}}
    \newcommand{\ControlFlowTok}[1]{\textcolor[rgb]{0.00,0.44,0.13}{\textbf{{#1}}}}
    \newcommand{\OperatorTok}[1]{\textcolor[rgb]{0.40,0.40,0.40}{{#1}}}
    \newcommand{\BuiltInTok}[1]{{#1}}
    \newcommand{\ExtensionTok}[1]{{#1}}
    \newcommand{\PreprocessorTok}[1]{\textcolor[rgb]{0.74,0.48,0.00}{{#1}}}
    \newcommand{\AttributeTok}[1]{\textcolor[rgb]{0.49,0.56,0.16}{{#1}}}
    \newcommand{\InformationTok}[1]{\textcolor[rgb]{0.38,0.63,0.69}{\textbf{\textit{{#1}}}}}
    \newcommand{\WarningTok}[1]{\textcolor[rgb]{0.38,0.63,0.69}{\textbf{\textit{{#1}}}}}
    
    
    % Define a nice break command that doesn't care if a line doesn't already
    % exist.
    \def\br{\hspace*{\fill} \\* }
    % Math Jax compatibility definitions
    \def\gt{>}
    \def\lt{<}
    \let\Oldtex\TeX
    \let\Oldlatex\LaTeX
    \renewcommand{\TeX}{\textrm{\Oldtex}}
    \renewcommand{\LaTeX}{\textrm{\Oldlatex}}
    % Document parameters
    % Document title
    \title{Numerical calculations of SMOx based gas sensors with Python}
    
    
    
    
    
% Pygments definitions
\makeatletter
\def\PY@reset{\let\PY@it=\relax \let\PY@bf=\relax%
    \let\PY@ul=\relax \let\PY@tc=\relax%
    \let\PY@bc=\relax \let\PY@ff=\relax}
\def\PY@tok#1{\csname PY@tok@#1\endcsname}
\def\PY@toks#1+{\ifx\relax#1\empty\else%
    \PY@tok{#1}\expandafter\PY@toks\fi}
\def\PY@do#1{\PY@bc{\PY@tc{\PY@ul{%
    \PY@it{\PY@bf{\PY@ff{#1}}}}}}}
\def\PY#1#2{\PY@reset\PY@toks#1+\relax+\PY@do{#2}}

\expandafter\def\csname PY@tok@w\endcsname{\def\PY@tc##1{\textcolor[rgb]{0.73,0.73,0.73}{##1}}}
\expandafter\def\csname PY@tok@c\endcsname{\let\PY@it=\textit\def\PY@tc##1{\textcolor[rgb]{0.25,0.50,0.50}{##1}}}
\expandafter\def\csname PY@tok@cp\endcsname{\def\PY@tc##1{\textcolor[rgb]{0.74,0.48,0.00}{##1}}}
\expandafter\def\csname PY@tok@k\endcsname{\let\PY@bf=\textbf\def\PY@tc##1{\textcolor[rgb]{0.00,0.50,0.00}{##1}}}
\expandafter\def\csname PY@tok@kp\endcsname{\def\PY@tc##1{\textcolor[rgb]{0.00,0.50,0.00}{##1}}}
\expandafter\def\csname PY@tok@kt\endcsname{\def\PY@tc##1{\textcolor[rgb]{0.69,0.00,0.25}{##1}}}
\expandafter\def\csname PY@tok@o\endcsname{\def\PY@tc##1{\textcolor[rgb]{0.40,0.40,0.40}{##1}}}
\expandafter\def\csname PY@tok@ow\endcsname{\let\PY@bf=\textbf\def\PY@tc##1{\textcolor[rgb]{0.67,0.13,1.00}{##1}}}
\expandafter\def\csname PY@tok@nb\endcsname{\def\PY@tc##1{\textcolor[rgb]{0.00,0.50,0.00}{##1}}}
\expandafter\def\csname PY@tok@nf\endcsname{\def\PY@tc##1{\textcolor[rgb]{0.00,0.00,1.00}{##1}}}
\expandafter\def\csname PY@tok@nc\endcsname{\let\PY@bf=\textbf\def\PY@tc##1{\textcolor[rgb]{0.00,0.00,1.00}{##1}}}
\expandafter\def\csname PY@tok@nn\endcsname{\let\PY@bf=\textbf\def\PY@tc##1{\textcolor[rgb]{0.00,0.00,1.00}{##1}}}
\expandafter\def\csname PY@tok@ne\endcsname{\let\PY@bf=\textbf\def\PY@tc##1{\textcolor[rgb]{0.82,0.25,0.23}{##1}}}
\expandafter\def\csname PY@tok@nv\endcsname{\def\PY@tc##1{\textcolor[rgb]{0.10,0.09,0.49}{##1}}}
\expandafter\def\csname PY@tok@no\endcsname{\def\PY@tc##1{\textcolor[rgb]{0.53,0.00,0.00}{##1}}}
\expandafter\def\csname PY@tok@nl\endcsname{\def\PY@tc##1{\textcolor[rgb]{0.63,0.63,0.00}{##1}}}
\expandafter\def\csname PY@tok@ni\endcsname{\let\PY@bf=\textbf\def\PY@tc##1{\textcolor[rgb]{0.60,0.60,0.60}{##1}}}
\expandafter\def\csname PY@tok@na\endcsname{\def\PY@tc##1{\textcolor[rgb]{0.49,0.56,0.16}{##1}}}
\expandafter\def\csname PY@tok@nt\endcsname{\let\PY@bf=\textbf\def\PY@tc##1{\textcolor[rgb]{0.00,0.50,0.00}{##1}}}
\expandafter\def\csname PY@tok@nd\endcsname{\def\PY@tc##1{\textcolor[rgb]{0.67,0.13,1.00}{##1}}}
\expandafter\def\csname PY@tok@s\endcsname{\def\PY@tc##1{\textcolor[rgb]{0.73,0.13,0.13}{##1}}}
\expandafter\def\csname PY@tok@sd\endcsname{\let\PY@it=\textit\def\PY@tc##1{\textcolor[rgb]{0.73,0.13,0.13}{##1}}}
\expandafter\def\csname PY@tok@si\endcsname{\let\PY@bf=\textbf\def\PY@tc##1{\textcolor[rgb]{0.73,0.40,0.53}{##1}}}
\expandafter\def\csname PY@tok@se\endcsname{\let\PY@bf=\textbf\def\PY@tc##1{\textcolor[rgb]{0.73,0.40,0.13}{##1}}}
\expandafter\def\csname PY@tok@sr\endcsname{\def\PY@tc##1{\textcolor[rgb]{0.73,0.40,0.53}{##1}}}
\expandafter\def\csname PY@tok@ss\endcsname{\def\PY@tc##1{\textcolor[rgb]{0.10,0.09,0.49}{##1}}}
\expandafter\def\csname PY@tok@sx\endcsname{\def\PY@tc##1{\textcolor[rgb]{0.00,0.50,0.00}{##1}}}
\expandafter\def\csname PY@tok@m\endcsname{\def\PY@tc##1{\textcolor[rgb]{0.40,0.40,0.40}{##1}}}
\expandafter\def\csname PY@tok@gh\endcsname{\let\PY@bf=\textbf\def\PY@tc##1{\textcolor[rgb]{0.00,0.00,0.50}{##1}}}
\expandafter\def\csname PY@tok@gu\endcsname{\let\PY@bf=\textbf\def\PY@tc##1{\textcolor[rgb]{0.50,0.00,0.50}{##1}}}
\expandafter\def\csname PY@tok@gd\endcsname{\def\PY@tc##1{\textcolor[rgb]{0.63,0.00,0.00}{##1}}}
\expandafter\def\csname PY@tok@gi\endcsname{\def\PY@tc##1{\textcolor[rgb]{0.00,0.63,0.00}{##1}}}
\expandafter\def\csname PY@tok@gr\endcsname{\def\PY@tc##1{\textcolor[rgb]{1.00,0.00,0.00}{##1}}}
\expandafter\def\csname PY@tok@ge\endcsname{\let\PY@it=\textit}
\expandafter\def\csname PY@tok@gs\endcsname{\let\PY@bf=\textbf}
\expandafter\def\csname PY@tok@gp\endcsname{\let\PY@bf=\textbf\def\PY@tc##1{\textcolor[rgb]{0.00,0.00,0.50}{##1}}}
\expandafter\def\csname PY@tok@go\endcsname{\def\PY@tc##1{\textcolor[rgb]{0.53,0.53,0.53}{##1}}}
\expandafter\def\csname PY@tok@gt\endcsname{\def\PY@tc##1{\textcolor[rgb]{0.00,0.27,0.87}{##1}}}
\expandafter\def\csname PY@tok@err\endcsname{\def\PY@bc##1{\setlength{\fboxsep}{0pt}\fcolorbox[rgb]{1.00,0.00,0.00}{1,1,1}{\strut ##1}}}
\expandafter\def\csname PY@tok@kc\endcsname{\let\PY@bf=\textbf\def\PY@tc##1{\textcolor[rgb]{0.00,0.50,0.00}{##1}}}
\expandafter\def\csname PY@tok@kd\endcsname{\let\PY@bf=\textbf\def\PY@tc##1{\textcolor[rgb]{0.00,0.50,0.00}{##1}}}
\expandafter\def\csname PY@tok@kn\endcsname{\let\PY@bf=\textbf\def\PY@tc##1{\textcolor[rgb]{0.00,0.50,0.00}{##1}}}
\expandafter\def\csname PY@tok@kr\endcsname{\let\PY@bf=\textbf\def\PY@tc##1{\textcolor[rgb]{0.00,0.50,0.00}{##1}}}
\expandafter\def\csname PY@tok@bp\endcsname{\def\PY@tc##1{\textcolor[rgb]{0.00,0.50,0.00}{##1}}}
\expandafter\def\csname PY@tok@fm\endcsname{\def\PY@tc##1{\textcolor[rgb]{0.00,0.00,1.00}{##1}}}
\expandafter\def\csname PY@tok@vc\endcsname{\def\PY@tc##1{\textcolor[rgb]{0.10,0.09,0.49}{##1}}}
\expandafter\def\csname PY@tok@vg\endcsname{\def\PY@tc##1{\textcolor[rgb]{0.10,0.09,0.49}{##1}}}
\expandafter\def\csname PY@tok@vi\endcsname{\def\PY@tc##1{\textcolor[rgb]{0.10,0.09,0.49}{##1}}}
\expandafter\def\csname PY@tok@vm\endcsname{\def\PY@tc##1{\textcolor[rgb]{0.10,0.09,0.49}{##1}}}
\expandafter\def\csname PY@tok@sa\endcsname{\def\PY@tc##1{\textcolor[rgb]{0.73,0.13,0.13}{##1}}}
\expandafter\def\csname PY@tok@sb\endcsname{\def\PY@tc##1{\textcolor[rgb]{0.73,0.13,0.13}{##1}}}
\expandafter\def\csname PY@tok@sc\endcsname{\def\PY@tc##1{\textcolor[rgb]{0.73,0.13,0.13}{##1}}}
\expandafter\def\csname PY@tok@dl\endcsname{\def\PY@tc##1{\textcolor[rgb]{0.73,0.13,0.13}{##1}}}
\expandafter\def\csname PY@tok@s2\endcsname{\def\PY@tc##1{\textcolor[rgb]{0.73,0.13,0.13}{##1}}}
\expandafter\def\csname PY@tok@sh\endcsname{\def\PY@tc##1{\textcolor[rgb]{0.73,0.13,0.13}{##1}}}
\expandafter\def\csname PY@tok@s1\endcsname{\def\PY@tc##1{\textcolor[rgb]{0.73,0.13,0.13}{##1}}}
\expandafter\def\csname PY@tok@mb\endcsname{\def\PY@tc##1{\textcolor[rgb]{0.40,0.40,0.40}{##1}}}
\expandafter\def\csname PY@tok@mf\endcsname{\def\PY@tc##1{\textcolor[rgb]{0.40,0.40,0.40}{##1}}}
\expandafter\def\csname PY@tok@mh\endcsname{\def\PY@tc##1{\textcolor[rgb]{0.40,0.40,0.40}{##1}}}
\expandafter\def\csname PY@tok@mi\endcsname{\def\PY@tc##1{\textcolor[rgb]{0.40,0.40,0.40}{##1}}}
\expandafter\def\csname PY@tok@il\endcsname{\def\PY@tc##1{\textcolor[rgb]{0.40,0.40,0.40}{##1}}}
\expandafter\def\csname PY@tok@mo\endcsname{\def\PY@tc##1{\textcolor[rgb]{0.40,0.40,0.40}{##1}}}
\expandafter\def\csname PY@tok@ch\endcsname{\let\PY@it=\textit\def\PY@tc##1{\textcolor[rgb]{0.25,0.50,0.50}{##1}}}
\expandafter\def\csname PY@tok@cm\endcsname{\let\PY@it=\textit\def\PY@tc##1{\textcolor[rgb]{0.25,0.50,0.50}{##1}}}
\expandafter\def\csname PY@tok@cpf\endcsname{\let\PY@it=\textit\def\PY@tc##1{\textcolor[rgb]{0.25,0.50,0.50}{##1}}}
\expandafter\def\csname PY@tok@c1\endcsname{\let\PY@it=\textit\def\PY@tc##1{\textcolor[rgb]{0.25,0.50,0.50}{##1}}}
\expandafter\def\csname PY@tok@cs\endcsname{\let\PY@it=\textit\def\PY@tc##1{\textcolor[rgb]{0.25,0.50,0.50}{##1}}}

\def\PYZbs{\char`\\}
\def\PYZus{\char`\_}
\def\PYZob{\char`\{}
\def\PYZcb{\char`\}}
\def\PYZca{\char`\^}
\def\PYZam{\char`\&}
\def\PYZlt{\char`\<}
\def\PYZgt{\char`\>}
\def\PYZsh{\char`\#}
\def\PYZpc{\char`\%}
\def\PYZdl{\char`\$}
\def\PYZhy{\char`\-}
\def\PYZsq{\char`\'}
\def\PYZdq{\char`\"}
\def\PYZti{\char`\~}
% for compatibility with earlier versions
\def\PYZat{@}
\def\PYZlb{[}
\def\PYZrb{]}
\makeatother


    % For linebreaks inside Verbatim environment from package fancyvrb. 
    \makeatletter
        \newbox\Wrappedcontinuationbox 
        \newbox\Wrappedvisiblespacebox 
        \newcommand*\Wrappedvisiblespace {\textcolor{red}{\textvisiblespace}} 
        \newcommand*\Wrappedcontinuationsymbol {\textcolor{red}{\llap{\tiny$\m@th\hookrightarrow$}}} 
        \newcommand*\Wrappedcontinuationindent {3ex } 
        \newcommand*\Wrappedafterbreak {\kern\Wrappedcontinuationindent\copy\Wrappedcontinuationbox} 
        % Take advantage of the already applied Pygments mark-up to insert 
        % potential linebreaks for TeX processing. 
        %        {, <, #, %, $, ' and ": go to next line. 
        %        _, }, ^, &, >, - and ~: stay at end of broken line. 
        % Use of \textquotesingle for straight quote. 
        \newcommand*\Wrappedbreaksatspecials {% 
            \def\PYGZus{\discretionary{\char`\_}{\Wrappedafterbreak}{\char`\_}}% 
            \def\PYGZob{\discretionary{}{\Wrappedafterbreak\char`\{}{\char`\{}}% 
            \def\PYGZcb{\discretionary{\char`\}}{\Wrappedafterbreak}{\char`\}}}% 
            \def\PYGZca{\discretionary{\char`\^}{\Wrappedafterbreak}{\char`\^}}% 
            \def\PYGZam{\discretionary{\char`\&}{\Wrappedafterbreak}{\char`\&}}% 
            \def\PYGZlt{\discretionary{}{\Wrappedafterbreak\char`\<}{\char`\<}}% 
            \def\PYGZgt{\discretionary{\char`\>}{\Wrappedafterbreak}{\char`\>}}% 
            \def\PYGZsh{\discretionary{}{\Wrappedafterbreak\char`\#}{\char`\#}}% 
            \def\PYGZpc{\discretionary{}{\Wrappedafterbreak\char`\%}{\char`\%}}% 
            \def\PYGZdl{\discretionary{}{\Wrappedafterbreak\char`\$}{\char`\$}}% 
            \def\PYGZhy{\discretionary{\char`\-}{\Wrappedafterbreak}{\char`\-}}% 
            \def\PYGZsq{\discretionary{}{\Wrappedafterbreak\textquotesingle}{\textquotesingle}}% 
            \def\PYGZdq{\discretionary{}{\Wrappedafterbreak\char`\"}{\char`\"}}% 
            \def\PYGZti{\discretionary{\char`\~}{\Wrappedafterbreak}{\char`\~}}% 
        } 
        % Some characters . , ; ? ! / are not pygmentized. 
        % This macro makes them "active" and they will insert potential linebreaks 
        \newcommand*\Wrappedbreaksatpunct {% 
            \lccode`\~`\.\lowercase{\def~}{\discretionary{\hbox{\char`\.}}{\Wrappedafterbreak}{\hbox{\char`\.}}}% 
            \lccode`\~`\,\lowercase{\def~}{\discretionary{\hbox{\char`\,}}{\Wrappedafterbreak}{\hbox{\char`\,}}}% 
            \lccode`\~`\;\lowercase{\def~}{\discretionary{\hbox{\char`\;}}{\Wrappedafterbreak}{\hbox{\char`\;}}}% 
            \lccode`\~`\:\lowercase{\def~}{\discretionary{\hbox{\char`\:}}{\Wrappedafterbreak}{\hbox{\char`\:}}}% 
            \lccode`\~`\?\lowercase{\def~}{\discretionary{\hbox{\char`\?}}{\Wrappedafterbreak}{\hbox{\char`\?}}}% 
            \lccode`\~`\!\lowercase{\def~}{\discretionary{\hbox{\char`\!}}{\Wrappedafterbreak}{\hbox{\char`\!}}}% 
            \lccode`\~`\/\lowercase{\def~}{\discretionary{\hbox{\char`\/}}{\Wrappedafterbreak}{\hbox{\char`\/}}}% 
            \catcode`\.\active
            \catcode`\,\active 
            \catcode`\;\active
            \catcode`\:\active
            \catcode`\?\active
            \catcode`\!\active
            \catcode`\/\active 
            \lccode`\~`\~ 	
        }
    \makeatother

    \let\OriginalVerbatim=\Verbatim
    \makeatletter
    \renewcommand{\Verbatim}[1][1]{%
        %\parskip\z@skip
        \sbox\Wrappedcontinuationbox {\Wrappedcontinuationsymbol}%
        \sbox\Wrappedvisiblespacebox {\FV@SetupFont\Wrappedvisiblespace}%
        \def\FancyVerbFormatLine ##1{\hsize\linewidth
            \vtop{\raggedright\hyphenpenalty\z@\exhyphenpenalty\z@
                \doublehyphendemerits\z@\finalhyphendemerits\z@
                \strut ##1\strut}%
        }%
        % If the linebreak is at a space, the latter will be displayed as visible
        % space at end of first line, and a continuation symbol starts next line.
        % Stretch/shrink are however usually zero for typewriter font.
        \def\FV@Space {%
            \nobreak\hskip\z@ plus\fontdimen3\font minus\fontdimen4\font
            \discretionary{\copy\Wrappedvisiblespacebox}{\Wrappedafterbreak}
            {\kern\fontdimen2\font}%
        }%
        
        % Allow breaks at special characters using \PYG... macros.
        \Wrappedbreaksatspecials
        % Breaks at punctuation characters . , ; ? ! and / need catcode=\active 	
        \OriginalVerbatim[#1,codes*=\Wrappedbreaksatpunct]%
    }
    \makeatother

    % Exact colors from NB
    \definecolor{incolor}{HTML}{303F9F}
    \definecolor{outcolor}{HTML}{D84315}
    \definecolor{cellborder}{HTML}{CFCFCF}
    \definecolor{cellbackground}{HTML}{F7F7F7}
    
    % prompt
    \makeatletter
    \newcommand{\boxspacing}{\kern\kvtcb@left@rule\kern\kvtcb@boxsep}
    \makeatother
    \newcommand{\prompt}[4]{
        \ttfamily\llap{{\color{#2}[#3]:\hspace{3pt}#4}}\vspace{-\baselineskip}
    }
    

    
    % Prevent overflowing lines due to hard-to-break entities
    \sloppy 
    % Setup hyperref package
    \hypersetup{
      breaklinks=true,  % so long urls are correctly broken across lines
      colorlinks=true,
      urlcolor=urlcolor,
      linkcolor=linkcolor,
      citecolor=citecolor,
      }
    % Slightly bigger margins than the latex defaults
    
    \geometry{verbose,tmargin=1in,bmargin=1in,lmargin=1in,rmargin=1in}
    
    

\begin{document}
    
    \maketitle
    
    

    
    \tableofcontents 
\setcounter{section}{0}

    \hypertarget{introduction}{%
\section{Introduction}\label{introduction}}

\hypertarget{motivation}{%
\subsection{Motivation}\label{motivation}}

When I first started working in the field of metal-oxide semiconducting
gas sensors, I still had the wish to finish my time at university with a
degree of eduction in math and physics. But while having spend 4 month
in a school and experiencing the joy of teaching, I also felt a strong
urge to continue learning and explore the possibilities science and
research would allow me. Long story short, I ended up starting my Phd.
in the ``Institute of Theoretical and Physical Chemistry'' at the
University of Tübingen. Lucky as I was I directly got the possibility of
working on a project with an industrial partner about gas sensors. This
allowed me to investigate the sate-of-the-art of industrialized sensor
production. Besides the benefits of learning to give regular
presentations and reports, working in teams and funding for interesting
research activities, it also imposed some new problems needed to be
solved. The biggest of challenges has been the continuous increase of
sensor date in always shorter intervals. Luck as I was, I was not the
only one facing problems with the increasing amount of data needed to be
analyzed. Aside from the conventional way of wiring
applications/programs, the need was rather on having efficient tool sets
to reduce time consuming parts of the data analysis. Especially those
which could theoretically be done by an algorithm. At this point my very
basic knowledge in the area of programming and the need to solve urgent
tasks, created the need for efficient solutions with a steep learning
curve. Until then my, and of some others in the lab, standard procedure
of working with our data was:

\begin{itemize}
\tightlist
\item
  first to extract relevant features at well defined time-points to
  reduce the initial amount to a reasonable quantity of data points
\item
  in a second step a unified representation of the extracted features
  had to be created
\item
  analysis of the data and taking further steps
\end{itemize}

When Working with a limited amount of samples, manually doing these
steps was acceptable. With the industrial cooperation gaining speed, the
number of samples increased also rapidly. Soon we reached the point
where the pre-processing of the data would easily consume a large amount
of time. And that without even having started with the analysis. Luckily
I was not the only one facing such problems and a project from the
Python community gained more and more interest. Articles like
\cite{Unpingco2014} pointed me to a new way of dealing with data and
using the Python programming language. Regarding the fact, that Python
was already a well estblished programming language, the introduction of
the ``IPtyhon notebook ecosystem'' was making the use of Python for in
an scientific workflow very attractive. As mentioned in this article:
\cite{Osipov2016}

\begin{quote}
``\ldots{} what they do offer is an environment for exploration,
collaboration, and visualization.''
\end{quote}

I also realized the large potential for my research field. Not only that
I got efficient tools for calculations, analysis and representations,
but the new tools have been build with the focus of easy reporting
results. Such an environment around the so called notebooks was the
ideal piece, which I experienced as a missing block in the scientific
work I was after. Besides my work for our industrial partner I also did
fundamental research about semiconducting metaloxide gas sensors. On the
basis of the great works on gas sensors, my focus was on one the
numerically calculation of the gas sensors semiconducting properties.
Typically a model of a gas sensor is developed and compared with
experimental results. The pees review system assures, that the
publications published in the journals are well written, documented the
on a solid and proven basis. Unfortunately the transfer from the
presented model in the publication to other model and a direct
comparison of the results experimental data aired on other sensors is
often not done. One of the main obstacles for the average experimental
oriented researcher is the lack of the programming knowledge to
implement algorithm based on the presented work.

With the results I gained while my Phd. thesis I had the experience,
that others understood the value of my work but could not transfer it to
their particular problem. At this point I decided to orient myself back
to my educational roots and use the powerful ecosystem around the
IPython notebooks and write my thesis with the large focus on suppling a
introduction to this excellent toolbox.

The presented results will allow more people to gain insides in the
semiconductors influence on the sensing properties. More than this I
hope this thesis will also be a useful introduction into Python,
specially IPython notebooks, for scientific work.

These hopes are not unfunded.In my time working at the ``Eberhard Karls
University of Tübingen'' i was also assigned to give multiple lectures.
One lecture branch has been the introduction to data mining. Often the
lack of programming knowledge and the limited amount of time did not
allow a usage of Python as a supporting programming language. In these
cases classical tools like Excel of Origin have been used to mine data
examples for hidden facts. Most attendees have been very fast and
understood the general concept of data mining. Just the lack of a the
minimal requirement to translate the concept into machine understandable
instructions stopped them from using them. When enough time was
available, a short introduction into programming with Python gained a
lot of interest and was generally seen as a positive experience. This
thesis is therefore structured in such a way, that one the one side I
will present my research results from the past years in a condensed form
and on the other side uses this the opportunity to introduce and explain
the importance of applying programming tools in the common work flow of
scientific work. Since the potential of ``outsourcing'' repetitive tasks
to machine executed scripts and invest more time in our ceativity and
intelligence is huge. My hope is to bring with this thesis not only a
deeper insight in the understanding of SMOX based gas sensors, but also
help others to start a interesting journey into the wide area of data
mining and machine learning with python.

I typically finish my introductions to Python by letting the students
run their first commands. Typically for other introductions found around
the globe, this is a program which outputs ``Hello World''. For Python I
prefer to execute some other commands with boils almost down to the
philosophical essence on how ``instructions'' should be.

    \begin{tcolorbox}[breakable, size=fbox, boxrule=1pt, pad at break*=1mm,colback=cellbackground, colframe=cellborder]
\prompt{In}{incolor}{1}{\boxspacing}
\begin{Verbatim}[commandchars=\\\{\}]
\PY{k+kn}{import} \PY{n+nn}{this}
\end{Verbatim}
\end{tcolorbox}

    \begin{Verbatim}[commandchars=\\\{\}]
The Zen of Python, by Tim Peters

Beautiful is better than ugly.
Explicit is better than implicit.
Simple is better than complex.
Complex is better than complicated.
Flat is better than nested.
Sparse is better than dense.
Readability counts.
Special cases aren't special enough to break the rules.
Although practicality beats purity.
Errors should never pass silently.
Unless explicitly silenced.
In the face of ambiguity, refuse the temptation to guess.
There should be one-- and preferably only one --obvious way to do it.
Although that way may not be obvious at first unless you're Dutch.
Now is better than never.
Although never is often better than *right* now.
If the implementation is hard to explain, it's a bad idea.
If the implementation is easy to explain, it may be a good idea.
Namespaces are one honking great idea -- let's do more of those!
    \end{Verbatim}

    \hypertarget{smox-based-gas-sensors}{%
\subsection{SMOX based gas sensors}\label{smox-based-gas-sensors}}

This section list a collection of high quality publications which will
broadcast all relevant information about semiconductor based gas sensors
directly into your brain! In the course of this thesis many of those
aspects will be explained, analyzed in discussed in required detail at
the relevant places. So reading these publications now is not necessary
to follow this thesis. This section is therefore a good point to come
back if a more detailed view on the subject is desired. Especially when
this thesis is used to apply the presented concepts to individual
research topics.

When solving numerical problems it typical to introduction some
assumptions, which are only valid under specific boundary conditions.
These assumptions simplify the problem to a level where a numerical
calculation is possible, but will reduce the validity of the results
only to a small subset of all possible situations. The calculations in
this thesis are also packed with assumptions and boundary conditions. My
intention is to supply enough information to understand the relevance of
the assumption and it's implications. With my presented work I do not
claim to calculate all aspects of SMOX gas sensors, but present a tool
which helps to understand specific aspects. The way of presenting this
knowledge should lead/motivate others to adapt the presented work to
individual other cases with completely different boundary conditions. A
general understanding of gas sensors will only be gained to a certain
level, but with time, work and data this understanding will definitely
grow.

    \begin{figure}
\centering
\includegraphics{media/icons/jupyter.png}
\caption{Juypter}
\end{figure}

    \hypertarget{notebooks}{%
\subsection{Notebooks}\label{notebooks}}

``The Zen of Python'' might not always be the primary directive of each
developer, but the Python community consists most probably of many
people how would consider the latter points as important. So did also
the inventors of the IPython and Jupyter. A quick search will reveal
multiple sources in the world wide web giving a detailed picture about
what Notebooks are and how Jupyter in connected with them. Here I will
not try to give an general overview about this tool and rather stick to
the phrase ``Learning by doing.''. By explaining topic related parts
this notebooks will guide an interested reader to the point where:

\begin{itemize}
\tightlist
\item
  understanding the fundamental instructions of Python
\item
  using the basic functionality of the Notebooks
\item
  fundamental understanding of SMOX based gas sensors
\end{itemize}

is gained.

It is worth mentioning, that the intention of such notebooks is to merge
the essential tools of scientific work flows together. Data acquisition,
preparation, analysis, representation and documentation should available
in one place. The strength of not just sharing final conclusions in a
nicely formatted way, but also being able to share the full stack of
steps necessary to reach the final conclusion is essentially the
strength of the Jupyter notebooks. This feature is already changing the
way how scientific results are shared/published and was intentionally
designed this way \cite{Randles2017}.

The formatting used in this notebook is based on Markdown. Wikipedia
summarizes Markdown like this:

\begin{quote}
Markdown is a lightweight markup language with plain text formatting
syntax. Wikipedia
\end{quote}

This means a document is formatted by writing plaintext and special text
blocks are interpreted as formatting commands. E.g. \textbf{BOLD}
letters are generated by encapsulating the text with
\texttt{**\ Text\ here\ **}, headings are generated by starting the
heading with \texttt{\#}. Depending on the number of \texttt{\#},
subsections are created. I will not go into detail here about the
features of Markdown. Many features are used in this notebook and are
directly accessible by double-clicking the text element. The plain text
will reveal the way it was created, and the execution of the cell with
\texttt{CTRL+ENTER} will reveal its Markdown formatted representation.
One other handy feature of the Jupyter ecosystem I use is the ability to
transform notebooks into multiple other formats. Just to name some:
HTML, WORD (DOC, DOCX), Latex. The tool \texttt{nbconvert} is used
internally to convert the Markdown formatted representation into other
formats. For this thesis the default option ``Export as PDF'' under the
File option generates a Latex based PDF file.
\href{https://guides.github.com/features/mastering-markdown/}{Mastering-markdown}
is a web page, where I found some hints on how to format my notebook.
For instance I gained the ability to make block quotes from this page
based on this example:

As Kanye West said:

\begin{quote}
We're living the future so the present is our past.
\end{quote}

    \hypertarget{installation-guide}{%
\subsubsection{Installation guide}\label{installation-guide}}

The easiest way to get started would be to use the Anaconda
distribution. Anaconda bundles multiple different tools and installs
them i nthe operation system. Anaconda will take care of cross
dependecies and handle the update process of the software. This is not
the only way to get started with Jupyter Notebooks but surely an very
fast and easy one.
\href{media/presentations/InstallAnaconda/Anaconda-Python.pdf}{HERE} is
additionally a presentation I use for my lectures to guide students into
the world of Python and
\href{media/presentations/InstallAnaconda/PythonÜbung.zip}{here one
example} of it's usage.

    \hypertarget{example-notebook---sneak-preview}{%
\subsubsection{Example Notebook - Sneak
preview}\label{example-notebook---sneak-preview}}

Besides the first example \texttt{import\ this}, here is a very basic
example which should prepare exited reader on whats coming next.

    \hypertarget{simple-is-better-than-complex.}{%
\paragraph{``Simple is better than
complex.''}\label{simple-is-better-than-complex.}}

The Jupyter environment is equipped with ``magic commands'', which are
not part of the base programming language (i.e.~Python), but rather a
helper instruction to simplify common tasks. Magic commands always start
with \texttt{\%} and are followed with an instruction. I will
demonstrate in this example the use of the \texttt{\%pylab} instruction.
This modifies the current programming space to be a lab nicely equipped
for scientific work flows. For instance a chemistry lab is equipped with
a balance, a water tab and a fire extinguisher, and in this case a
``pylab'' is equipped with data handing, plotting and multiple other
tools. So let's setup a ``pylab'' and run some lines of code. (The
plotting is handled in the background by Matplotlib \cite{Hunter2007})

    \begin{tcolorbox}[breakable, size=fbox, boxrule=1pt, pad at break*=1mm,colback=cellbackground, colframe=cellborder]
\prompt{In}{incolor}{2}{\boxspacing}
\begin{Verbatim}[commandchars=\\\{\}]
\PY{o}{\PYZpc{}}\PY{k}{pylab} inline

\PY{c+c1}{\PYZsh{}get a list of 5000 points between \PYZhy{}50 and 50}
\PY{n}{x} \PY{o}{=} \PY{n}{linspace}\PY{p}{(}\PY{n}{start}\PY{o}{=}\PY{o}{\PYZhy{}}\PY{l+m+mi}{50}\PY{p}{,} \PY{n}{stop}\PY{o}{=}\PY{l+m+mi}{50}\PY{p}{,} \PY{n}{num}\PY{o}{=}\PY{l+m+mi}{5000}\PY{p}{)}

\PY{c+c1}{\PYZsh{} Python way of adding a comment is the \PYZsq{}\PYZsh{}\PYZsq{} character}
\PY{c+c1}{\PYZsh{} Python way of writer \PYZsq{}x to the power of y\PYZsq{} is x**y}
\PY{c+c1}{\PYZsh{} here I am calculating thethe third power of x for 5000 points between \PYZhy{}50 and 50}
\PY{n}{y} \PY{o}{=} \PY{n}{x}\PY{o}{*}\PY{o}{*}\PY{l+m+mi}{3}

\PY{c+c1}{\PYZsh{}plot it}
\PY{n}{fig} \PY{o}{=} \PY{n}{plt}\PY{o}{.}\PY{n}{plot}\PY{p}{(}\PY{n}{x}\PY{p}{,}\PY{n}{y}\PY{p}{,}\PY{l+s+s1}{\PYZsq{}}\PY{l+s+s1}{b}\PY{l+s+s1}{\PYZsq{}}\PY{p}{)}
\end{Verbatim}
\end{tcolorbox}

    \begin{Verbatim}[commandchars=\\\{\}]
Populating the interactive namespace from numpy and matplotlib
    \end{Verbatim}

    \begin{center}
    \adjustimage{max size={0.9\linewidth}{0.9\paperheight}}{1_Introduction_files/1_Introduction_9_1.png}
    \end{center}
    { \hspace*{\fill} \\}
    
    The additional ``inline'' instruction in \texttt{\%pylab} configures the
plotting environment to represent the graphics inside this text. The
following line generates a 5000 linear distributed points in the
interval {[}-50,50) and y is just the third power of each point.
\texttt{plt.} is a submodule defined by the magic command
\texttt{\%pylab} which handels data plotting in a very simple way.
\texttt{plt.plot(x,y,\textquotesingle{}r-\textquotesingle{})} for
example plots x vs.~y with a red line. The thesis is done in the
notebook to offer the reader the possibility to directly work with newly
gained knowledge. Therefore the upper block is a good opportunity to do
the fist steps in Python. For instance change the line format to `b'
(blue). Or to `b-.'. Change the exponent from 3 to 3.1. To do so:

\begin{itemize}
\tightlist
\item
  click on the code box
\item
  edit the field
\item
  Add the following line
  \texttt{plt.title(\textquotesingle{}The\ power\ law\textquotesingle{})}
\item
  go back with CTRL-Z to correct your mistakes
\item
  hit CTRL-ENTER to execute the code
\end{itemize}

    \hypertarget{conclusion}{%
\section{Conclusion}\label{conclusion}}

Since the motivation for this ``interactive'' thesis should be clear
now, I would like to come in the next section now to my actual research
topic: ``Numerical calculation of semiconducting metal oxide (SMOX)
based gas sensors''. In the \href{2-SMOX-sensor-calculation.ipynb}{next
chapter} I will demonstrate how theoretical numerical calculations for
chemical sensors are used to better understand experimental results. In
the \href{3-Applying_theorie-to-practice.ipynb}{following chapter} I
will demonstrate how such knowledge could be used to improve the
performance of a sensor regarding it's sensitivity and selectivity.

\href{2-SMOX-sensor-calculation.ipynb}{Follow this link to come to the
next section.}

    \hypertarget{about-the-pdf-version-of-this-work}{%
\section{About the PDF-Version of this
work}\label{about-the-pdf-version-of-this-work}}

This notebook was not intended to be used as a printed hard copy or as a
PDF. The provided PDF servers just as an low level representation of the
original work. It should give potentially interested readers an easy
entry point to Python(any other programming language) supported science.
Many of the implemented features in these notebooks like interactive
widgets, animated data representations and live code examples will not
work in the PDF-version. Only a static snapshot of the mutable
representation can be represented outside the notebook. On the one side
the reduction of printed pages and the benefits of interacting with the
presented work in a notebook should motivate the reader to use the
notebook. Nevertheless the integrated function to export the notebook to
a Latex based PDF results in a very nice static publication. So please
keep in mind, that the PDF-verison may not be able to represent all the
features of the notebook as intended and some links might not work as
expected. You are strongly encouraged to switch to the Jupyter
presentation of this work and experience the full potential of such
notebooks.

\hypertarget{equations}{%
\subsection{Equations}\label{equations}}

In a typical scientific thesis and textbook, relevant equations are
referred by numbers of identifiers. In Latex this is done, by assigning
a label and a tag to an equation. If this equation needs to be referred
to, the a pointer to the reference is added, and the tag is used for the
representation of this equation. As an example, an arbitrary equation
from this thesis is used. The (internal) reference of this equation is
`second\_derivative', while the printed representation (tag) is ``Second
derivative''.

\begin{align}
\frac{d²V^{*}}{dr^{*2}}=1-n^{*}(V^{*})-\frac{2}{r^{*}}\frac{dV^{*}}{dr^{*}}\label{second_derivative}\tag{Second derivative}
\end{align}

To refer not to this equation a reference can be added which will look
like this (\ref{second_derivative}). The underlying mechanisms will link
the reference with the equation, print the label, and add an hyperlink,
to be able to jump to the equation. This feature works either in
notebook or PDF representations. One drawback of the notebook
representation is, that he equation and reference need to be in the same
code block/cell. Otherwise the reference is not working. Instead of
linking the equation correctly, ??? will be represented. My opinion is,
that this will change will evolving improvements of Jupyter is is just a
temporal flaw. Since the PDF version is processed by a Latex interpreter
in total, the ``one code block'' limitation does not exist there. To
demonstrate this, I will try to reference the eqaution in the next code
block.

    Here the same reference to the equation: \ref{second_derivative}

    \hypertarget{tables}{%
\subsection{Tables}\label{tables}}

With the numerical calculations performed in this thesis, data will need
to be represented also. One way of doing this is by plotting the the
data. This is also suitable for a static PDF-Version of this thesis,
which often is useful/necessary to have. This feature is well
implemented and the transfer from Notebook to PDF version works well.
Another way to display data are tables. Similar to the well known,
omnipresent tool for working with tables, Excel, the Python universe has
its own, but similar, tool. It's called Pandas. As a primer I will give
here a very short introduction to Pandas. What sheets are for Excel, are
Dataframes for Pandas. Here a simple example how to create a Dataframe
in analogy to the previous programming example:

    \begin{tcolorbox}[breakable, size=fbox, boxrule=1pt, pad at break*=1mm,colback=cellbackground, colframe=cellborder]
\prompt{In}{incolor}{3}{\boxspacing}
\begin{Verbatim}[commandchars=\\\{\}]
\PY{o}{\PYZpc{}}\PY{k}{pylab}
\PY{k+kn}{import} \PY{n+nn}{pandas}

\PY{c+c1}{\PYZsh{}create some x values}
\PY{n}{x} \PY{o}{=} \PY{n}{linspace}\PY{p}{(}\PY{n}{start}\PY{o}{=}\PY{l+m+mi}{0}\PY{p}{,}\PY{n}{stop}\PY{o}{=}\PY{l+m+mi}{5}\PY{p}{,}\PY{n}{num} \PY{o}{=} \PY{l+m+mi}{10}\PY{p}{)}

\PY{c+c1}{\PYZsh{}the y values will the the square of the x values}
\PY{n}{y} \PY{o}{=} \PY{n}{x}\PY{o}{*}\PY{o}{*}\PY{l+m+mi}{2}

\PY{c+c1}{\PYZsh{}Put them in a Dataframe and reference this new Dataframe with the variable `dF`}
\PY{n}{dF} \PY{o}{=} \PY{n}{pandas}\PY{o}{.}\PY{n}{DataFrame}\PY{p}{(}\PY{p}{\PYZob{}}\PY{l+s+s1}{\PYZsq{}}\PY{l+s+s1}{x}\PY{l+s+s1}{\PYZsq{}}\PY{p}{:}\PY{n}{x}\PY{p}{,} \PY{l+s+s1}{\PYZsq{}}\PY{l+s+s1}{y}\PY{l+s+s1}{\PYZsq{}}\PY{p}{:}\PY{n}{y}\PY{p}{\PYZcb{}}\PY{p}{)}
\end{Verbatim}
\end{tcolorbox}

    \begin{Verbatim}[commandchars=\\\{\}]
Using matplotlib backend: Qt5Agg
Populating the interactive namespace from numpy and matplotlib
    \end{Verbatim}

    This simple example of a Dataframe should demonstrate the basic concept
of Dataframes. Once data is in the \texttt{DataFrame} format, the are
infinit ways to transform/slice/merge/\ldots{} it, to bring it into the
desired shape. Along this thesis, some of the functionalities of
Dataframes will be used and explained. Since the Jupyter environment is
still ``under construction'', the transfer from Juypter to PDF for
notebooks does not work to well until now. This does not mean, that is
is not possible, it is just not yet implemented in the default/vanilla
environment. This is way I want to highlight the small tweak that is at
this point still needed to have a comparable output in Juypter notebooks
and printed PDFs. To bypass this obstacle the default behavior of pandas
needs to be altered ( similar to this post:
\href{https://stackoverflow.com/a/24167756}{Latex-Tables Monkey Patch}).
This following patch brings DataFrames in the appropriate shape to by
nicely represented in both representations.

    \hypertarget{before-the-patch}{%
\subsubsection{Before the patch}\label{before-the-patch}}

    \begin{tcolorbox}[breakable, size=fbox, boxrule=1pt, pad at break*=1mm,colback=cellbackground, colframe=cellborder]
\prompt{In}{incolor}{4}{\boxspacing}
\begin{Verbatim}[commandchars=\\\{\}]
\PY{n}{display}\PY{p}{(}\PY{n}{dF}\PY{p}{)}
\end{Verbatim}
\end{tcolorbox}

    
    \begin{verbatim}
          x          y
0  0.000000   0.000000
1  0.555556   0.308642
2  1.111111   1.234568
3  1.666667   2.777778
4  2.222222   4.938272
5  2.777778   7.716049
6  3.333333  11.111111
7  3.888889  15.123457
8  4.444444  19.753086
9  5.000000  25.000000
    \end{verbatim}

    
    \hypertarget{the-patch}{%
\subsubsection{The patch}\label{the-patch}}

    \begin{tcolorbox}[breakable, size=fbox, boxrule=1pt, pad at break*=1mm,colback=cellbackground, colframe=cellborder]
\prompt{In}{incolor}{5}{\boxspacing}
\begin{Verbatim}[commandchars=\\\{\}]
\PY{k+kn}{import} \PY{n+nn}{pandas}
\PY{n}{pandas}\PY{o}{.}\PY{n}{set\PYZus{}option}\PY{p}{(}\PY{l+s+s1}{\PYZsq{}}\PY{l+s+s1}{display.notebook\PYZus{}repr\PYZus{}html}\PY{l+s+s1}{\PYZsq{}}\PY{p}{,} \PY{k+kc}{True}\PY{p}{)}

\PY{n}{pandas}\PY{o}{.}\PY{n}{set\PYZus{}option}\PY{p}{(}\PY{l+s+s1}{\PYZsq{}}\PY{l+s+s1}{display.notebook\PYZus{}repr\PYZus{}html}\PY{l+s+s1}{\PYZsq{}}\PY{p}{,} \PY{k+kc}{True}\PY{p}{)}

\PY{k}{def} \PY{n+nf}{\PYZus{}repr\PYZus{}latex\PYZus{}}\PY{p}{(}\PY{n+nb+bp}{self}\PY{p}{)}\PY{p}{:}
    \PY{k}{return} \PY{l+s+sa}{r}\PY{l+s+s2}{\PYZdq{}\PYZdq{}\PYZdq{}}
\PY{l+s+s2}{    }\PY{l+s+s2}{\PYZbs{}}\PY{l+s+s2}{begin}\PY{l+s+si}{\PYZob{}center\PYZcb{}}
\PY{l+s+s2}{    }\PY{l+s+s2}{\PYZob{}}\PY{l+s+si}{\PYZpc{}s}\PY{l+s+s2}{\PYZcb{}}
\PY{l+s+s2}{    }\PY{l+s+s2}{\PYZbs{}}\PY{l+s+s2}{end}\PY{l+s+si}{\PYZob{}center\PYZcb{}}
\PY{l+s+s2}{    }\PY{l+s+s2}{\PYZdq{}\PYZdq{}\PYZdq{}} \PY{o}{\PYZpc{}} \PY{n+nb+bp}{self}\PY{o}{.}\PY{n}{to\PYZus{}latex}\PY{p}{(}\PY{p}{)}

\PY{n}{pandas}\PY{o}{.}\PY{n}{DataFrame}\PY{o}{.}\PY{n}{\PYZus{}repr\PYZus{}latex\PYZus{}} \PY{o}{=} \PY{n}{\PYZus{}repr\PYZus{}latex\PYZus{}}  \PY{c+c1}{\PYZsh{} monkey patch pandas DataFrame}
\end{Verbatim}
\end{tcolorbox}

    \hypertarget{prettier-tables-after-the-patch}{%
\subsubsection{Prettier tables after the
patch}\label{prettier-tables-after-the-patch}}

    \begin{tcolorbox}[breakable, size=fbox, boxrule=1pt, pad at break*=1mm,colback=cellbackground, colframe=cellborder]
\prompt{In}{incolor}{6}{\boxspacing}
\begin{Verbatim}[commandchars=\\\{\}]
\PY{n}{display}\PY{p}{(}\PY{n}{dF}\PY{p}{)}
\PY{c+c1}{\PYZsh{}\PYZdq{}Test Tabel\PYZdq{}}
\end{Verbatim}
\end{tcolorbox}

    
    \begin{center}
    {\begin{tabular}{lrr}
\toprule
{} &         x &          y \\
\midrule
0 &  0.000000 &   0.000000 \\
1 &  0.555556 &   0.308642 \\
2 &  1.111111 &   1.234568 \\
3 &  1.666667 &   2.777778 \\
4 &  2.222222 &   4.938272 \\
5 &  2.777778 &   7.716049 \\
6 &  3.333333 &  11.111111 \\
7 &  3.888889 &  15.123457 \\
8 &  4.444444 &  19.753086 \\
9 &  5.000000 &  25.000000 \\
\bottomrule
\end{tabular}
}
    \end{center}
    

    
    In Jupyter notebooks the output will look very similar. In the PDF
version of this thesis, those two output will differ(, unless a newer
version of Jupyter fixed this issue.) To have this thesis in a PDF
printable form, I will use the presented patch to unify the output. On
the other side I suggest not to use this patch, and rather work with the
Jupyter notebooks and its default representation, since the patch has
also its downsides. But this would go beyond the scope of this remark
about the patch.

    \hypertarget{bibliography-section}{%
\section{Bibliography section}\label{bibliography-section}}

    \bibliographystyle{alphadin}
\bibliography{ipython}


    % Add a bibliography block to the postdoc
    
    
    
\end{document}
